% cSpell:locale en,pt-BR - Adicionado para suporte ao corretor ortográfico
% Extensão Code Spell Checker no VS Code (muito útil!)

\documentclass[12pt,a4paper,oneside]{book}
% O TIPO DE DOCUMENTO é opção deste pacote e governa o modelo inteiro:
% natureza, folha de aprovação, resumos e metadados exigidos. Seis valores:
%
%   generico     relatório de disciplina (padrão)      NBR 14724 no que couber
%   estagio      relatório de estágio                  NBR 14724 "e similares"
%   tcc1         projeto de pesquisa                   NBR 15287:2025
%   tcc2         monografia, grau Bacharel             NBR 14724:2024
%   dissertacao  monografia, grau Mestre               NBR 14724:2024
%   tese         monografia, grau Doutor               NBR 14724:2024
%
% Troque o valor abaixo para mudar de tipo. Sem a opção, vale generico.
% O tipo vai como CHAVE (tipo=...), e não solto: assim um valor digitado
% errado dá erro deste pacote, nomeando os aceitos, em vez de virar erro de
% idioma do babel — que é para onde vão as opções que o pacote não reconhece.
\usepackage[tipo=generico]{UnifeiICTReport}
%\usepackage[utf8]{inputenc} % Removido pois XeLaTeX usa UTF-8 nativo
\usepackage{csquotes}
\usepackage{fontspec}
\usepackage{import}
\usepackage{lipsum} % Para gerar texto placeholder
\usepackage[style=abnt]{biblatex}
\usepackage{fancyvrb} % Para formatação avançada de verbatim
\usepackage{subcaption}
\usepackage{booktabs}

% \usepackage{showframe} % para debug de margens e layout
% \usepackage{lua-visual-debug} % para debug visual com LuaLaTeX

\makeglossaries % necessário para glossários e listas de abreviaturas
\subimport{Preambulo/}{lista-abreviaturas.tex}
\subimport{Preambulo/}{lista-simbolos.tex}

\bibliography{referencias.bib}

%--------------------------------------
% Informações de dados para CAPA e FOLHA DE ROSTO
%--------------------------------------

% Título, subtítulo e autores (compatibilidade com unifeiICT.sty)
\title{Modelo de Relatório Técnico-Científico para os cursos do ICT}

% Comente ou remova a linha abaixo se não houver subtítulo.
% Escreva SÓ o texto do subtítulo, sem dois-pontos: a NBR 14724:2024 §4.1.1(d)
% exige que ele venha "precedido de dois-pontos, evidenciando a sua
% subordinação ao título", e o modelo insere essa pontuação sozinho na capa,
% na folha de rosto e na folha de aprovação.
\subtitle{Incluindo instruções de uso}

% Se houver mais de um autor, separe os nomes com \\ (quebra de linha)
\author{Autor do Relatório\\Outro autor do Relatório\\Mais um autor do Relatório}

% Se o relatório for orientado por um professor, descomente a linha abaixo e preencha o nome.
% Geralmente, o orientador é o professor da disciplina para a qual o relatório está sendo elaborado.
\supervisor{Prof. Dr. Orientador Exemplar}

% Descomente a linha abaixo e preencha o nome se houver coorientador
% \cosupervisor{Prof. Dra. Coorientadora Exemplar}

% Disciplina que o relatório atende
\subject{Nome da Disciplina}

% Metadados adicionais usados pela folha de aprovação (\folhaaprovacao, mais
% abaixo). Data de aprovação e nota são opcionais; sem elas, a folha sai sem
% essas linhas, em vez de com campos em branco estranhos.
% \aprovacaodata{15 de dezembro de 2025}
% \notaaprovacao{9,5}

% Só para os tipos de monografia (tcc2, dissertacao, tese).
% A área de concentração é EXIGIDA em dissertacao e tese — sem ela a
% compilação para, nomeando o que falta. No tcc2 ela é opcional e rara: só
% quando o trabalho integra um projeto maior. A linha de pesquisa é sempre
% opcional.
% \areaconcentracao{Sistemas de Computação}
% \linhapesquisa{Engenharia de Software}

% Membros da banca examinadora — usados por tcc1 (quando há defesa), tcc2,
% dissertacao e tese. Declare apenas os membros EXTERNOS: o orientador já vem
% de \supervisor e entra na banca sozinho, sem precisar ser repetido aqui.
% Campos em branco ainda rendem um bloco corretamente formatado, para a folha
% preparada antes de os nomes serem conhecidos.
% No tcc1 a defesa é facultativa: sem nenhum \bancamembro declarado, a folha
% de aprovação simplesmente não é emitida.
% \bancamembro{Prof. Dr. Fulano de Tal}{Doutor em Engenharia}{Unifei}
% \bancamembro{Prof. Dra. Beltrana de Tal}{Doutora em Engenharia}{Unifei}


% Resumo na língua principal
\abstract{%
    Este é um modelo de resumo para o relatório técnico-científico. O resumo deve apresentar de forma  concisa os objetivos, a metodologia, os resultados e as conclusões do trabalho. Deve ser escrito em um único parágrafo, com no máximo 250 palavras, e deve ser claro e objetivo, permitindo ao leitor compreender rapidamente o conteúdo do relatório. O resumo principal deve ser escrito na língua principal do relatório.
}
% Palavras-chave do resumo principal
\keywords{relatório técnico-científico; modelo; Unifei ICT; LaTeX.}

% Resumo na língua secundária, geralmente inglês.
\abstractseclang{%
    This is a model abstract for the technical-scientific report. The abstract should concisely present the objectives, methodology, results, and conclusions of the work. It should be
    written in a single paragraph, with a maximum of 250 words, and should be clear and objective,
    allowing the reader to quickly understand the content of the report. The secondary abstract should  be written in English. If the report is written in English, the secondary abstract should be in Portuguese.
}
% Palavras-chave do resumo secundário
\keywordsseclang{technical-scientific report; model; Unifei ICT; LaTeX.}

% ---
% Configurações de pacotes
% ---

\begin{document}

% Preâmbulo do documento (elementos pré-textuais)
\frontmatter

% ---
% Impressão da capa e folha de rosto
% ---

\maketitle

% ---
% Folha de aprovação (NBR 14724:2024 §4.2.1.3) — sem argumento: a forma dela
% vem do tipo declarado lá em cima, no \usepackage.
% No tcc1 ela só sai se houver banca declarada; nos demais tipos é obrigatória.
% ---
\folhaaprovacao

% ---
% Resumos
% ---
\makeabstracts

%----
%  Listas de figuras, tabelas, abreviaturas e símbolos
%  -- Estes elementos são opcionais e podem ser incluídos conforme a necessidade do relatório.
%  -- Para inclui-los, descomente as linhas abaixo. Para suprimir, deixe comentado.
% ----

\listoffigures

\listoftables

% Quadro e Gráfico (NBR 14724:2024 §5.8) têm lista própria, como Figura e
% Tabela. Descomente a que fizer sentido para o relatório; nenhum tipo usado
% no texto gera página em branco por a lista ficar comentada.
% \listofquadros
%
% \listofgraficos

\printglossary[type=\acronymtype,title={\acronymlistname}]

\printglossary[type=symbols,title={\symbolslistname}]

% ---
% Sumário
% ---
\tableofcontents

\mainmatter

% ---
% Introdução
% ---
\subimport{Capitulos/cap1/}{cap1.tex}

% ---
% Desenvolvimento
% ---
\subimport{Capitulos/cap2/}{cap2.tex}

% ---
% Conclusão
% ---
\subimport{Capitulos/cap3/}{cap3.tex}

% ---
% Referências bibliográficas
% --- Imprime a bibliografia
% ---
\printbibliography

% ---
% Elementos pós-textuais são iniciados após \backmatter
% --- Compreende apêndices, anexos, glossários e índices
% --- São todos opcionais
% ---

\backmatter

% Apêndices: produzidos pelo próprio autor, para complementar a argumentação
% (NBR 14724:2024 §3.4). \apendices abre o grupo (uma entrada coletiva no
% sumário); cada \apendice{<rótulo>}{<título>} abre um apêndice, lettered A, B, ...
\apendices

\apendice{ape:questionario}{Questionário aplicado na pesquisa}
Este apêndice é de autoria do próprio autor do relatório — daí ser um
\emph{apêndice}, e não um anexo. Ilustrações dentro de um apêndice continuam a
mesma sequência numérica do restante do documento e aparecem normalmente nas
listas, como a \reffig{fig:questionario} a seguir demonstra.

\begin{figuraabnt}{fig:questionario}{Estrutura do questionário aplicado}
    \includegraphics[width=0.6\textwidth]{example-image}
    \fonte{o próprio autor}
\end{figuraabnt}

% Anexos: documentos de terceiros, incorporados como fundamentação, prova ou
% ilustração (§3.3) -- não escritos pelo autor do relatório.
\anexos

\anexo{ane:certificado}{Certificado de participação no evento}
Este anexo reproduz um documento de terceiros — o certificado emitido pela
organização do evento — anexado ao relatório como prova, e não redigido pelo
autor. Essa autoria é exatamente o que separa um \emph{anexo} de um
\emph{apêndice}: veja o \refcomp{}{ape:questionario}, de autoria própria, em
contraste com este.

\end{document}